\BLOCK{include '_preamble.tex'}

\begin{document}
\maketitle

%% ---- POCETAK PRORACUNA ----

\subsection{Ulazni parametri}

\begin{table}[H]
\centering
\begin{tabular}{|c|l|c|c|}
\hline
Oznaka & Opis & Vrednost & Jedinica \\
\hline
$N_{Ed}$ & Normalna sila pritiska & \VAR{m.NEd|n} & kN \\
\hline
$b$ & Širina preseka & \VAR{(m.b*100)|n(1)} & cm \\
\hline
$h$ & Visina preseka & \VAR{(m.h*100)|n(1)} & cm \\
\hline
$f_{ck}$ & Karakteristična čvrstoća betona & \VAR{(m.fck/1000)|n(0)} & MPa \\
\hline
$q_0$ & Osnovna vrednost faktora ponašanja & \VAR{m.q0|n} & - \\
\hline
$O$ & Obim uzengija & \VAR{(m.ObimUzg*100)|n(1)} & cm \\
\hline
- & Tip armature & \VAR{m.tip_armature} & - \\
\hline
$\phi_w$ & Prečnik uzengija & \VAR{precnik} & mm \\
\hline
$T$ & Period oscilovanja & \VAR{m.T|n} & s \\
\hline
\BLOCK{set pola = (bi|length + 1) // 2}
\BLOCK{if bi|length > 1}
%# Razmaci b_i u dva reda tabele (bez linije između), da stanu u širinu stranice
$b_i$ & Razmaci pridržanih šipki & \VAR{bi[:pola]|map('n', 1)|join(', ')}, & cm \\
 & & \VAR{bi[pola:]|map('n', 1)|join(', ')} & \\
\BLOCK{else}
$b_i$ & Razmaci pridržanih šipki & \VAR{bi|map('n', 1)|join(', ')} & cm \\
\BLOCK{endif}
\hline
\end{tabular}
\caption{Ulazni parametri stuba \VAR{naziv|tex}}
\end{table}

\begin{figure}[H]
\centering
\includegraphics[width=0.55\textwidth]{\VAR{slike.presek}}
\caption{Poprečni presek stuba sa oznakama}
\end{figure}

\subsection{Materijal i geometrija utegnutog jezgra}

\begin{align*}
f_{cd} &= \frac{0.85 \cdot f_{ck}}{1.5} = \frac{0.85 \cdot \VAR{(m.fck/1000)|n(0)}}{1.5} = \VAR{(m.fcd/1000)|n} \ \mathrm{MPa} \\
f_{yd} &= \frac{f_{yk}}{\gamma_s} = \frac{500}{1.15} = \VAR{(m.fyd/1000)|n} \ \mathrm{MPa} \\
\varepsilon_{sy,d} &= \VAR{(m.esyd*1000)|n(3)} \ \text{‰} \\
b_0 &= b - 8 = \VAR{(m.b*100)|n(1)} - 8 = \VAR{(m.b0*100)|n(1)} \ \mathrm{cm} \\
h_0 &= h - 8 = \VAR{(m.h*100)|n(1)} - 8 = \VAR{(m.h0*100)|n(1)} \ \mathrm{cm} \\
A_{sw} &= \frac{\phi_w^2 \cdot \pi}{4} = \VAR{(m.izabrani_fi*10000)|n(3)} \ \mathrm{cm^2}
\end{align*}

\subsection{Normalizovana aksijalna sila}

\begin{align*}
\nu_{Ed} &= \frac{N_{Ed}}{b \cdot h \cdot f_{cd}} = \frac{\VAR{m.NEd|n}}{\VAR{m.b|n(3)} \cdot \VAR{m.h|n(3)} \cdot \VAR{(m.fcd/1000)|n} \cdot 10^3} = \VAR{m.nied|n(3)}
\end{align*}

Kontrola maksimalne normalizovane sile:

\begin{align*}
\nu_{Ed} &\leq 0.65 \\
\VAR{m.nied|n(3)} &\leq 0.65 \quad \VAR{kontrola(m.nied <= 0.65)}
\end{align*}

\BLOCK{if d.slucaj == 'prekoracena_sila'}
\textbf{Zaključak:} normalizovana sila je veća od dozvoljene. Potrebno je povećati marku betona ili dimenzije stuba; dalji proračun utezanja se ne sprovodi.
\BLOCK{elif d.slucaj == 'EC2'}
Kako je $\nu_{Ed} = \VAR{m.nied|n(3)} < 0.2$ i $q_0 = \VAR{m.q0|n} \leq 2$ (EN 1998-1, 5.4.3.2.2(12)P):

\textbf{Zaključak:} utezanje se ne proračunava, primenjuju se konstruktivna pravila iz EC2.
\BLOCK{else}

\subsection{Faktor duktilnosti krivine $\mu_\phi$}

Prema EN 1998-1, 5.4.3.2.2(6)P i 5.2.3.4(3), sa osnovnom vrednošću faktora ponašanja $q_0$.
Za armaturu \VAR{m.tip_armature} koeficijent prema 5.2.3.4(4) iznosi $k = \VAR{k_arm|n(1)}$.

\BLOCK{if m.T >= m.Tc}
Kako je $T_1 = \VAR{m.T|n} \ \mathrm{s} \geq T_C = \VAR{m.Tc|n} \ \mathrm{s}$, izraz (5.4):

\begin{align*}
\mu_\phi &= k \cdot (2 \cdot q_0 - 1) = \VAR{k_arm|n(1)} \cdot (2 \cdot \VAR{m.q0|n} - 1) = \VAR{d.mfi|n(3)}
\end{align*}
\BLOCK{else}
Kako je $T_1 = \VAR{m.T|n} \ \mathrm{s} < T_C = \VAR{m.Tc|n} \ \mathrm{s}$, izraz (5.5):

\begin{align*}
\mu_\phi &= k \cdot \left( 1 + 2 \cdot (q_0 - 1) \cdot \frac{T_C}{T_1} \right) = \VAR{k_arm|n(1)} \cdot \left( 1 + 2 \cdot (\VAR{m.q0|n} - 1) \cdot \frac{\VAR{m.Tc|n}}{\VAR{m.T|n}} \right) = \VAR{d.mfi|n(3)}
\end{align*}
\BLOCK{endif}

\subsection{Koeficijent efikasnosti utezanja u preseku $\alpha_n$}

\begin{align*}
\alpha_n &= 1 - \frac{\sum b_i^2}{6 \cdot b_0 \cdot h_0} = 1 - \frac{\VAR{bi|map('n', 1)|join('^2 + ')}^2}{6 \cdot \VAR{(m.b0*100)|n(1)} \cdot \VAR{(m.h0*100)|n(1)}} = \VAR{d.alfan|n(4)}
\end{align*}

\subsection{Mehanički zapreminski procenat utezanja}

Zapremina uzengija po jednom redu:

\begin{align*}
V_{sw} &= A_{sw} \cdot O = \VAR{(m.izabrani_fi*10000)|n(3)} \cdot \VAR{(m.ObimUzg*100)|n(1)} = \VAR{(d.Vsw*1e6)|n(2)} \ \mathrm{cm^3}
\end{align*}

Za svaki razmak uzengija $s$ (redom 20, 15, 12.5, 10 i 7.5 cm) proverava se uslov $\omega_{wd,prov} \geq \omega_{wd,req}$, gde je:

\begin{align*}
\alpha_s &= \left(1 - \frac{s}{2 \cdot b_0}\right) \cdot \left(1 - \frac{s}{2 \cdot h_0}\right), \qquad \alpha = \alpha_n \cdot \alpha_s \\
\omega_{wd,prov} &= \frac{V_{sw} \cdot f_{yd}}{b_0 \cdot h_0 \cdot s \cdot f_{cd}} \\
\alpha \cdot \omega_{wd,req} &= 30 \cdot \mu_\phi \cdot \nu_{Ed} \cdot \varepsilon_{sy,d} \cdot \frac{b}{b_0} - 0.035 = 30 \cdot \VAR{d.mfi|n(3)} \cdot \VAR{m.nied|n(3)} \cdot \VAR{m.esyd|n(6)} \cdot \frac{\VAR{(m.b*100)|n(1)}}{\VAR{(m.b0*100)|n(1)}} - 0.035 = \VAR{iteracije[0].alfa_omegawdreq|n(4)} \\
\omega_{wd,req} &= \max\left(0.08;\ \frac{\alpha \cdot \omega_{wd,req}}{\alpha}\right) \quad \text{(EN 1998-1, 5.4.3.2.2(9))}
\end{align*}

\begin{table}[H]
\centering
\begin{tabular}{|c|c|c|c|c|c|}
\hline
$s$ [cm] & $\alpha_s$ & $\alpha$ & $\omega_{wd,prov}$ & $\omega_{wd,req}$ & Kontrola \\
\hline
\BLOCK{for it in iteracije}
\VAR{(it.s*100)|n(1)} & \VAR{it.alfas|n(4)} & \VAR{it.alfa|n(4)} & \VAR{it.omegawdprov|n(4)} & \VAR{it.omegawdreq|n(4)} & \VAR{'zadovoljava' if it.zadovoljava else 'ne zadovoljava'} \\
\hline
\BLOCK{endfor}
\end{tabular}
\caption{Provera razmaka uzengija $\phi$\VAR{precnik}}
\end{table}

\BLOCK{set it = iteracije[-1]}
\BLOCK{if d.uspeh}
Usvojeni razmak uzengija je $s = \VAR{(it.s*100)|n(1)} \ \mathrm{cm}$:
\BLOCK{else}
Ni najmanji razmak ne zadovoljava; prikaz za $s = \VAR{(it.s*100)|n(1)} \ \mathrm{cm}$:
\BLOCK{endif}

\begin{align*}
\alpha_s &= \left(1 - \frac{\VAR{(it.s*100)|n(1)}}{2 \cdot \VAR{(m.b0*100)|n(1)}}\right) \cdot \left(1 - \frac{\VAR{(it.s*100)|n(1)}}{2 \cdot \VAR{(m.h0*100)|n(1)}}\right) = \VAR{it.alfas|n(4)} \\
\alpha &= \alpha_n \cdot \alpha_s = \VAR{d.alfan|n(4)} \cdot \VAR{it.alfas|n(4)} = \VAR{it.alfa|n(4)} \\
\omega_{wd,prov} &= \frac{\VAR{(d.Vsw*1e6)|n(2)} \cdot \VAR{(m.fyd/1000)|n}}{\VAR{(m.b0*100)|n(1)} \cdot \VAR{(m.h0*100)|n(1)} \cdot \VAR{(it.s*100)|n(1)} \cdot \VAR{(m.fcd/1000)|n}} = \VAR{it.omegawdprov|n(4)} \\
\omega_{wd,req} &= \max\left(0.08;\ \frac{\VAR{it.alfa_omegawdreq|n(4)}}{\VAR{it.alfa|n(4)}}\right) = \VAR{it.omegawdreq|n(4)}
\end{align*}

Kontrola utezanja:

\begin{align*}
\omega_{wd,prov} &\geq \omega_{wd,req} \\
\VAR{it.omegawdprov|n(4)} &\geq \VAR{it.omegawdreq|n(4)} \quad \VAR{kontrola(it.zadovoljava)}
\end{align*}

\subsection{Zaključak}

\BLOCK{if d.uspeh}
Stub \VAR{naziv|tex} je dovoljno utegnut uzengijama $\phi$\VAR{precnik} na razmaku $s = \VAR{(it.s*100)|n(1)} \ \mathrm{cm}$.
\BLOCK{else}
\textbf{Nedovoljno utezanje stuba \VAR{naziv|tex} za izabrani prečnik uzengija $\phi$\VAR{precnik}.} Potrebno je povećati prečnik uzengija ili obim (broj krakova) uzengija.
\BLOCK{endif}
\BLOCK{endif}

%% ---- KRAJ PRORACUNA ----

\end{document}
